\chapter{Overview}
\label{chap:overview}

% archon:covers MR4717077CanonicalRepresentationsOfSurfaceGroups.lean MR4717077CanonicalRepresentationsOfSurfaceGroups/Basic.lean

This chapter is the comprehensive blueprint for the formalization of
Landesman--Litt, \emph{Canonical representations of surface groups}
(MR4717077, arXiv:2205.15352).
The central result is \cref{thm:finite-image}: any MCG-finite representation
$\rho\colon \pi_1(\Sigma_{g,n}) \to \GL_r(\mathbb{C})$ with $r < \sqrt{g+1}$
has finite image.
The proof has three steps: (1)~the unitary case
(\cref{prop:unitary-case}), using cohomological vanishing and integrality of
cohomologically rigid local systems; (2)~the semisimple case
(\cref{thm:semisimple-case}), reducing to unitary via non-abelian Hodge
theory; (3)~the general case, using \cref{lem:semisimplicity} to reduce to
the semisimple case via the asymptotic Putman--Wieland theorem
(\cref{thm:putman-wieland-asymptotic}).
Throughout, the key analytic engine is the cohomological rank bound
(\cref{thm:rank-bound-subsystem}), proved via period map analysis
(\cref{thm:period-map}).

% The section structure below mirrors the section structure (titles and order)
% of the reference paper Landesman--Litt, "Canonical representations of surface
% groups" (references/MR4717077-surface-groups.tex). Items are grouped under the
% paper section they originate from (per their SOURCE provenance comments).

\section{Introduction}

% ---------------------------------------------------------------------------
% Definitions
% ---------------------------------------------------------------------------

\begin{definition}[Surface group]
  \label{def:surface-group}
  \lean{MR4717077CanonicalRepresentationsOfSurfaceGroups.SurfaceGroup}
  Fix non-negative integers $g, n$ with $\Sigma_{g,n}$ hyperbolic, i.e.\
  $n \geq 1$ if $g = 1$ and $n \geq 3$ if $g = 0$.
  The \emph{surface group} $\pi_1(\Sigma_{g,n})$ is the fundamental group of
  an orientable topological surface of genus $g$ with $n$ punctures.
  For $n = 0$ it has the standard presentation
  \[
    \langle a_1, b_1, \ldots, a_g, b_g \mid
      [a_1,b_1]\cdots[a_g,b_g] = 1 \rangle.
  \]
  For $n \geq 1$ the group $\pi_1(\Sigma_{g,n})$ is free on $2g + n - 1$
  generators.
\end{definition}

\begin{definition}[Mapping class group]
  \label{def:mapping-class-group}
  \lean{MR4717077CanonicalRepresentationsOfSurfaceGroups.MappingClassGroup}
  % SOURCE: [Landesman--Litt], §1.1 notation, (read from references/MR4717077-surface-groups.tex)
  % SOURCE QUOTE: "We will use $\on{Mod}_{g,n}$ to denote the mapping class group of an orientable surface of genus $g$ with $n$ punctures, and $\on{PMod}_{g,n}=\pi_1(\mathscr{M}_{g,n})$ (see \autoref{lemma:mgn-fundamental-group}) to denote pure mapping class group, i.e.~the subgroup of $\on{Mod}_{g,n}$ of index $n!$ preserving the punctures pointwise."
  \textit{Source: Landesman--Litt, §1.1 notation.}
  The \emph{mapping class group} $\Mod_{g,n} = \pi_0(\mathrm{Homeo}^+(\Sigma_{g,n}))$
  is the group of isotopy classes of orientation-preserving homeomorphisms of
  $\Sigma_{g,n}$ fixing the set of punctures setwise.
  It acts on $\pi_1(\Sigma_{g,n})$ via an outer action induced by the
  homeomorphism action.
  The \emph{pure mapping class group}
  $\PMod_{g,n} = \pi_1(\mathcal{M}_{g,n})$
  is the index-$n!$ subgroup of $\Mod_{g,n}$ preserving each puncture
  pointwise, where $\mathcal{M}_{g,n}$ is the Deligne--Mumford moduli stack
  of $n$-pointed genus-$g$ smooth proper curves.
\end{definition}

\begin{definition}[MCG-finite representation]
  \label{def:mcg-finite}
  \lean{MR4717077CanonicalRepresentationsOfSurfaceGroups.MCGFinite}
  \uses{def:surface-group, def:mapping-class-group}
  % SOURCE: [Landesman--Litt], Definition 1.1.1 (definition:mcg-finite), (read from references/MR4717077-surface-groups.tex)
  % SOURCE QUOTE: "For $g,n,r \geq 0$, a representation $\rho: \pi_1(\Sigma_{g,n},x) \to \on{GL}_r(\mathbb C)$ is {\em MCG-finite} if the conjugacy class of $\rho$ has finite orbit under the action of $\on{Mod}_{g,n}$."
  \textit{Source: Landesman--Litt, Definition 1.1.1.}
  For $g, n, r \geq 0$, a representation
  $\rho\colon \pi_1(\Sigma_{g,n}, x) \to \GL_r(\mathbb{C})$ is
  \emph{MCG-finite} if its conjugacy class has finite orbit under the action
  of $\Mod_{g,n}$ on conjugacy classes of representations.
\end{definition}

% ---------------------------------------------------------------------------
% Corollaries stated in the introduction (§1.2, §1.3, §1.6)
% ---------------------------------------------------------------------------

\begin{corollary}[Isomonodromy application]
  \label{cor:isomonodromy}
  \lean{MR4717077CanonicalRepresentationsOfSurfaceGroups.isomonodromy_iff_finiteMonodromy}
  \uses{thm:finite-image}
  % SOURCE: [Landesman--Litt], Corollary 1.2.3, (read from references/MR4717077-surface-groups.tex)
  % SOURCE QUOTE: "Let $C$ be a smooth projective curve of genus $g$, and let $D\subset C$ be a reduced effective divisor. Let $({E},\nabla)$ be a semisimple flat vector bundle on $C$ with regular singularities along $D$, with $\rk {E}<\sqrt{g+1}$. Suppose that the eigenvalues of the residue matrices of $({E},\nabla)$ have real parts in $[0,1)$. Then $({E},\nabla)$ has an algebraic universal isomonodromic deformation if and only if $({E},\nabla)$ has finite monodromy."
  \textit{Source: Landesman--Litt, Corollary 1.2.3.}
  Let $C$ be a smooth projective curve of genus $g$, $D \subset C$ a reduced
  effective divisor, and $(E,\nabla)$ a semisimple flat vector bundle on $C$
  with regular singularities along $D$, $\rk E < \sqrt{g+1}$, and eigenvalues
  of the residue matrices with real parts in $[0,1)$.
  Then $(E,\nabla)$ has an algebraic universal isomonodromic deformation if
  and only if $(E,\nabla)$ has finite monodromy.
\end{corollary}

\begin{proof}
  \uses{thm:finite-image}
  By Cousin--Heu (Theorem~A of \cite{CH:isomonodromic}), algebraic
  isomonodromic deformations correspond exactly to MCG-finite semisimple
  monodromy representations.
  Apply \cref{thm:finite-image} to deduce finite monodromy; the converse is
  immediate since finite monodromy gives a finite (hence MCG-finite) conjugacy
  class.
\end{proof}

\begin{corollary}[Local systems on versal families]
  \label{cor:versal-families}
  \lean{MR4717077CanonicalRepresentationsOfSurfaceGroups.finiteMonodromy_versalFamily}
  \uses{thm:finite-image, prop:mcg-finite-family}
  % SOURCE: [Landesman--Litt], Corollary 1.3.1, (read from references/MR4717077-surface-groups.tex)
  % SOURCE QUOTE: "Suppose $\pi^\circ: \mathscr C^\circ \to \mathscr M$ is a punctured versal family of genus $g$ curves, and $\mathbb V$ is a local system on $\mathscr C^\circ$ of rank $< \sqrt{g+1}$. For any fiber $C^\circ$ of $\pi^\circ$, the local system $\mathbb V|_{C^\circ}$ has finite monodromy."
  \textit{Source: Landesman--Litt, Corollary 1.3.1.}
  Let $\pi^\circ\colon \mathcal{C}^\circ \to \mathcal{M}$ be a punctured versal
  family of genus-$g$ curves, and $\mathbb{V}$ a local system on
  $\mathcal{C}^\circ$ of rank $r < \sqrt{g+1}$.
  For any fiber $C^\circ$ of $\pi^\circ$, the restriction $\mathbb{V}|_{C^\circ}$
  has finite monodromy.
\end{corollary}

\begin{proof}
  \uses{prop:mcg-finite-family, thm:finite-image}
  By \cref{prop:mcg-finite-family}, $\mathbb{V}|_{C^\circ}$ is MCG-finite.
  Apply \cref{thm:finite-image} with $r < \sqrt{g+1}$.
\end{proof}

\begin{corollary}[Free groups]
  \label{cor:free-groups}
  \lean{MR4717077CanonicalRepresentationsOfSurfaceGroups.freeGroup_finiteImage}
  \uses{thm:finite-image, def:surface-group}
  % SOURCE: [Landesman--Litt], Corollary 1.6.1 (corollary:free-groups), (read from references/MR4717077-surface-groups.tex)
  % SOURCE QUOTE: "Let $F_N$ be a free group on $N$ generators, with $N=2g$ or $N=2g+1$. Let $$\rho: F_N\to \on{GL}_r(\mathbb{C})$$ be a representation whose conjugacy class has finite orbit under $\on{Out}(F_N)$. If $r<\sqrt{g+1}$, then $\rho$ has finite image."
  \textit{Source: Landesman--Litt, Corollary 1.6.1.}
  Let $F_N$ be a free group on $N$ generators with $N = 2g$ or $N = 2g+1$.
  Let $\rho\colon F_N \to \GL_r(\mathbb{C})$ be a representation whose
  conjugacy class has finite orbit under $\Out(F_N)$.
  If $r < \sqrt{g+1}$, then $\rho$ has finite image.
\end{corollary}

% SOURCE QUOTE PROOF: "If $N$ is even, choose an isomorphism $F_N\simeq \pi_1(\Sigma_{g,1})$ and if $N$ is odd choose $F_N\simeq \pi_1(\Sigma_{g,2})$. Any representation of $F_N$ with finite orbit under $\on{Out}(F_N)$ yields a representation of a genus $g$ surface group with finite orbit under the mapping class group. Such a representation has finite image by \autoref{theorem:finite-image}."
\begin{proof}
  \uses{thm:finite-image, def:surface-group}
  For $N = 2g$, choose an isomorphism $F_N \cong \pi_1(\Sigma_{g,1})$; for
  $N = 2g+1$, choose $F_N \cong \pi_1(\Sigma_{g,2})$ (see
  \cref{def:surface-group}).
  Any representation of $F_N$ with finite orbit under $\Out(F_N)$ yields,
  under this identification, an MCG-finite representation of a genus-$g$
  surface group.
  Apply \cref{thm:finite-image}.
\end{proof}

\section{Representation-theoretic preliminaries}

% ---------------------------------------------------------------------------
% Group-theoretic preliminaries
% ---------------------------------------------------------------------------

\begin{lemma}[Birman exact sequence]
  \label{lem:birman-exact-sequence}
  \lean{MR4717077CanonicalRepresentationsOfSurfaceGroups.BirmanExactSequence}
  \uses{def:mapping-class-group, def:surface-group}
  % SOURCE: [Landesman--Litt], Lemma 2.0.1 (lemma:mgn-fundamental-group), (read from references/MR4717077-surface-groups.tex)
  % SOURCE QUOTE: "For $m \in \mathscr M_{g,n}$ a basepoint and $m'\in \mathscr M_{g, n+1}$ a lift of $m$, there are isomorphisms $\pi_1(\mathscr M_{g,n}, m) \simeq \on{PMod}_{g,n}$, $\pi_1(\mathscr M_{g,n+1}, m') \simeq \on{PMod}_{g,n}$ such that the diagram commutes, where the vertical maps are given by these isomorphisms, the top row is the exact sequence of homotopy groups induced by the forgetful map $\mathscr{M}_{g,n+1}\to \mathscr{M}_{g,n}$, and the bottom row is the Birman exact sequence \cite[p.~98]{farbM:a-primer}."
  \textit{Source: Landesman--Litt, Lemma 2.0.1.}
  For basepoints $m \in \mathcal{M}_{g,n}$ and $m' \in \mathcal{M}_{g,n+1}$
  lifting $m$, there are isomorphisms
  $\pi_1(\mathcal{M}_{g,n+1}, m') \cong \PMod_{g,n+1}$ and
  $\pi_1(\mathcal{M}_{g,n}, m) \cong \PMod_{g,n}$
  fitting into the Birman exact sequence
  \[
    1 \to \pi_1(\Sigma_{g,n}) \to \PMod_{g,n+1} \to \PMod_{g,n} \to 1,
  \]
  where the inclusion $\pi_1(\Sigma_{g,n}) \hookrightarrow \PMod_{g,n+1}$
  is the point-pushing subgroup.
\end{lemma}

\begin{proof}
  \uses{def:mapping-class-group, def:surface-group}
  % SOURCE QUOTE PROOF: "This follows from the contractibility of the universal cover of $\mathscr{M}_{g, n}$, and the fact that $\mathscr{M}_{g,n}$ is the quotient of its universal cover by the properly discontinuous action of $\on{PMod}_{g,n}$. See \cite[\S10.6.3 and p. 353]{farbM:a-primer}."
  This follows from the contractibility of the universal cover of
  $\mathcal{M}_{g,n}$ and the fact that $\mathcal{M}_{g,n}$ is the quotient
  of its universal cover by the properly discontinuous action of $\PMod_{g,n}$;
  see Farb--Margalit \S10.6.3 and p.~353.
\end{proof}

\begin{proposition}[Basic properties of MCG-finite representations]
  \label{prop:mcg-finite-basic}
  \lean{MR4717077CanonicalRepresentationsOfSurfaceGroups.mcgFinite_directSum}
  \uses{def:mcg-finite}
  % SOURCE: [Landesman--Litt], Proposition 2.1.1 (proposition:basic-properties), (read from references/MR4717077-surface-groups.tex)
  % SOURCE QUOTE: "The direct sum and tensor product of two MCG-finite representations of $\pi_1(\Sigma_{g,n})$ is MCG-finite. Any semisimple subquotient of an MCG-finite representation is MCG-finite."
  \textit{Source: Landesman--Litt, Proposition 2.1.1.}
  The direct sum and tensor product of two MCG-finite representations of
  $\pi_1(\Sigma_{g,n})$ is MCG-finite.
  Any semisimple subquotient of an MCG-finite representation is MCG-finite.
\end{proposition}

\begin{proof}
  \uses{def:mcg-finite}
  % SOURCE QUOTE PROOF: "The first statement, about direct sums and tensor products, is clear. The second is immediate from \cite[Lemma 2.2.1]{lawrence2019representations}."
  The claim about direct sums and tensor products is clear from the definition.
  The claim about semisimple subquotients follows since the stabilizer of any
  irreducible subquotient has finite index in the stabilizer of the original
  representation (Lawrence, Lemma 2.2.1).
\end{proof}

\begin{proposition}[MCG-finite from family of curves]
  \label{prop:mcg-finite-family}
  \lean{MR4717077CanonicalRepresentationsOfSurfaceGroups.mcgFinite_of_localSystem}
  \uses{def:mcg-finite, def:mapping-class-group, lem:birman-exact-sequence}
  % SOURCE: [Landesman--Litt], Proposition 2.1.3 (proposition:MCG-finite-family-of-curves), (read from references/MR4717077-surface-groups.tex)
  % SOURCE QUOTE: "Let $\pi^\circ: \mathscr{C}^\circ\to \mathscr{M}$ be a punctured versal family of $n$-pointed genus $g$ curves, as in \autoref{notation:versal-family}. In particular, $\mathscr M \to \mathscr M_{g,n}$ is dominant and \'etale. Let $C^\circ$ be a fiber of $\pi^\circ$. If $$\rho:\pi_1(\mathscr{C}^\circ)\to \on{GL}_r(\mathbb{C})$$ is a representation, $\rho|_{\pi_1(C^\circ)}$ is MCG-finite."
  \textit{Source: Landesman--Litt, Proposition 2.1.3.}
  Let $\pi^\circ\colon \mathcal{C}^\circ \to \mathcal{M}$ be a punctured versal
  family of $n$-pointed genus-$g$ curves, i.e.\ the map
  $\mathcal{M} \to \mathcal{M}_{g,n}$ is dominant \'{e}tale.
  If $\mathbb{V}$ is a local system on $\mathcal{C}^\circ$ and $C^\circ$ is any
  fiber of $\pi^\circ$, then $\mathbb{V}|_{C^\circ}$ is MCG-finite.
\end{proposition}

\begin{proof}
  \uses{def:mcg-finite, lem:birman-exact-sequence}
  The fibration sequence for $\mathcal{C}^\circ \to \mathcal{M}$ induces an
  outer action of $\pi_1(\mathcal{M})$ on $\pi_1(C^\circ)$, which factors
  through the outer action of $\PMod_{g,n}$ on $\pi_1(C^\circ)$ from the
  Birman exact sequence (\cref{lem:birman-exact-sequence}).
  Since $\mathcal{M} \to \mathcal{M}_{g,n}$ is dominant \'{e}tale, the image
  of $\pi_1(\mathcal{M}) \to \PMod_{g,n}$ has finite index.
  Any element of $\pi_1(\mathcal{C}^\circ)$ above such a class conjugates the
  restriction, preserving its conjugacy class; hence the stabilizer of the
  conjugacy class of $\mathbb{V}|_{C^\circ}$ has finite index in $\Mod_{g,n}$.
\end{proof}

\begin{lemma}[Tensor product decomposition for unitary local systems]
  \label{lem:tensor-product-unitary}
  \lean{MR4717077CanonicalRepresentationsOfSurfaceGroups.tensorDecomp_unitary}
  \uses{def:mcg-finite, prop:mcg-finite-basic, prop:mcg-finite-family}
  % SOURCE: [Landesman--Litt], Lemma 2.4.2 (lemma:unitary-direct-sum), (read from references/MR4717077-surface-groups.tex)
  % SOURCE QUOTE: "Let $A$ be an Artin local $\mathbb{C}$-algebra with residue field $\mathbb{C}$. Suppose $\mathbb V$ is a local system of free $A$-modules on $\mathscr C^\circ$ such that $\mathbb V|_{\mathscr C^\circ_m}$ is a constant deformation of a unitary local system, i.e.~there exists a unitary $\mathbb{C}$-local system $\mathbb{V}_0$ on $\mathscr C^\circ_m$ such that $\mathbb V|_{\mathscr C^\circ_m}\simeq \mathbb{V}_0\otimes_\mathbb{C} A$. There is a dominant \'etale map $\mathscr M' \to \mathscr M$, with $\mathscr{C}'=\mathscr{M}'\times_\mathscr{M}\mathscr{C}$ and ${\pi'}^\circ: {\mathscr{C}'}^\circ\to \mathscr{M}'$ the associated family of punctured curves, over which $\mathbb V|_{ {\mathscr{C}'}^\circ} \simeq \oplus_{i=1}^s \mathbb U_i \otimes ({\pi'}^\circ)^* \mathbb W_i$, where the $\mathbb W_i$ are locally constant sheaves of free $A$-modules on $\mathscr M'$ and the $\mathbb U_i$ are unitary local systems on ${\mathscr{C}'}^\circ$."
  \textit{Source: Landesman--Litt, Lemma 2.4.2.}
  Let $\mathbb{V}$ be a local system on $\mathcal{C}^\circ$ whose restriction
  to a fiber $C^\circ$ is unitary.
  After replacing $\mathcal{M}$ by a dominant \'{e}tale cover $\mathcal{M}'$,
  there is a decomposition
  \[
    \mathbb{V}|_{\mathcal{C}'^\circ}
    \cong \bigoplus_{i=1}^s \mathbb{U}_i \otimes (\pi'^\circ)^* \mathbb{W}_i,
  \]
  where each $\mathbb{U}_i$ is an irreducible unitary local system on
  $\mathcal{C}'^\circ$ and each $\mathbb{W}_i$ is a local system on
  $\mathcal{M}'$.
\end{lemma}

\begin{proof}
  \uses{def:mcg-finite, prop:mcg-finite-basic, prop:mcg-finite-family}
  Decompose the unitary fibral monodromy as a direct sum of irreducibles
  $\mathbb{V}_0 \cong \bigoplus_i \mathbb{S}_i^{n_i}$.
  Each $\mathbb{S}_i$ is MCG-finite by \cref{prop:mcg-finite-basic,prop:mcg-finite-family};
  lift each to a unitary local system $\mathbb{U}_i$ on a suitable finite
  \'{e}tale cover.
  Set $\mathbb{W}_i = \pi'^\circ_* \mathrm{Hom}(\mathbb{U}_i, \mathbb{V})$;
  the evaluation map $\bigoplus_i (\pi'^\circ)^* \mathbb{W}_i \otimes \mathbb{U}_i
  \to \mathbb{V}$ is an isomorphism fiberwise by Schur's lemma.
\end{proof}

\section{Hodge-theoretic preliminaries}

% ---------------------------------------------------------------------------
% Non-abelian Hodge theory input (§4.3 Consequences of non-abelian Hodge theory)
% ---------------------------------------------------------------------------

\begin{theorem}[Non-abelian Hodge unitarity]
  \label{thm:nah-unitarity}
  \lean{MR4717077CanonicalRepresentationsOfSurfaceGroups.nah_unitarity}
  \uses{thm:cohomologically-rigid-integral}
  % SOURCE: [Landesman--Litt 2022], Theorem 1.2.12 ([LL22]), cited in §1.9.1 and §1.9.2 of MR4717077, (read from references/MR4717077-surface-groups.tex)
  % SOURCE QUOTE: "these $\rho\otimes_{\mathcal O_{K,\iota}}\mathbb{C}$ underlie complex polarizable variations of Hodge structure for any complex structure on $\Sigma_{g,n}$. Hence, given their low rank, these local systems are unitary by \cite[Theorem 1.2.12]{LL22}."
  \textit{Source: Landesman--Litt [LL22], Theorem 1.2.12 (cited in MR4717077, §§1.9.1–1.9.2).}
  Let $\mathbb{V}$ be a semisimple, cohomologically rigid local system of rank
  $r < \sqrt{g+1}$ on a curve $C^\circ$ that underlies a complex polarizable
  variation of Hodge structure (i.e.\ is defined over $\mathcal{O}_K$ and each
  Galois conjugate is a VHS).
  Then $\mathbb{V}$ is unitary: its monodromy representation factors through
  the unitary group $U(r)$.
  (Axiom: relies on non-abelian Hodge theory developed in [LL22]; not yet
  available in Mathlib.)
\end{theorem}

\begin{proof}
  Proved in Landesman--Litt (arXiv:2202.05190), Theorem 1.2.12, using the
  Corlette--Simpson correspondence and the assumption that the local system is
  defined over a number ring with all Galois conjugates giving VHS.
  Since every complex polarizable VHS of rank $r < \sqrt{g+1}$ on a curve has
  unitary monodromy, the integrality from \cref{thm:cohomologically-rigid-integral}
  combined with the VHS property forces unitarity.
  Introduced as an axiom here pending formalization of non-abelian Hodge theory.
\end{proof}

\section{The period map associated to a unitary representation}

% ---------------------------------------------------------------------------
% Period map (§5.1) and bilinear pairing (§5.2)
% ---------------------------------------------------------------------------

\begin{theorem}[Period map derivative]
  \label{thm:period-map}
  \lean{MR4717077CanonicalRepresentationsOfSurfaceGroups.periodMap_derivative}
  % SOURCE: [Landesman--Litt], Theorem 5.1.6 (theorem:period-map-computation), (read from references/MR4717077-surface-groups.tex)
  % SOURCE QUOTE: "Under the identifications above, the map $dP^\vee_b$ factors as [diagram with $H^0(C, {E}\otimes \Omega^1_C(\log D))\otimes H^0(C, {E}^\vee \otimes \omega_C) \ar[r]^-{dP^\vee_b} \ar[d]^{\otimes} & \Omega^1_{\mathscr{B},b}$ and $H^0(C, {E}\otimes {E}^\vee \otimes \omega_C^{\otimes 2}(D)) \ar[r]^-{\on{tr}} & H^0(C, \omega_C^{\otimes 2}(D)) \ar[u]^{c^*_b}$] where the left vertical map $\otimes$ is the tensor product of global sections, and the bottom map $\on{tr}$ is induced by the trace pairing ${E}\otimes {E}^\vee\to \mathscr{O}_C$. Moreover, $dP^\vee_b$ is adjoint to $\overline \nabla_b$."
  \textit{Source: Landesman--Litt, Theorem 5.1.6 (proved in Appendix~A).}
  Let $C$ be a smooth projective curve with boundary divisor $D$, and let
  $E$ be the Mehta--Seshadri vector bundle on $C$ associated to a unitary
  local system on $C \setminus D$.
  The derivative of the period map associated to the mixed Hodge structure on
  $R^1\pi^\circ_* \mathbb{V}$ is identified with the natural multiplication map
  \[
    H^0(E \otimes \omega_C(D)) \otimes H^0(E^\vee \otimes \omega_C)
    \xrightarrow{\otimes}
    H^0(E \otimes E^\vee \otimes \omega_C^{\otimes 2}(D))
    \xrightarrow{\mathrm{tr}}
    H^0(\omega_C^{\otimes 2}(D)),
  \]
  where $\mathrm{tr}$ is induced by the trace pairing
  $E \otimes E^\vee \to \mathcal{O}_C$.
  (Introduced as an axiom pending Hodge theory infrastructure in Mathlib.)
\end{theorem}

\begin{proof}
  Proved in Appendix~A of Landesman--Litt (arXiv:2205.15352) via an explicit
  computation of the Kodaira--Spencer class and its interaction with the
  Mehta--Seshadri correspondence.  Introduced as an axiom in the Lean
  formalization, as the required mixed Hodge theory is not yet available in
  Mathlib.
\end{proof}

\begin{proposition}[Bilinear pairing rank bound]
  \label{proposition:pairing-result}
  \lean{MR4717077CanonicalRepresentationsOfSurfaceGroups.pairingResult}
  \uses{thm:period-map}
  % SOURCE: [Landesman--Litt], §5.2 (proposition:pairing-result), (read from references/MR4717077-surface-groups.tex)
  % SOURCE QUOTE: "Let $E_\star$ be a semistable parabolic bundle on $(C,D)$ of parabolic degree zero, with underlying vector bundle $E$. Let $v\in H^0(C, E\otimes \omega_C(D))$ be a nonzero global section, and suppose the map $H^0(C, B_E(v,-)): H^0(C, E^\vee \otimes \omega_C) \rightarrow H^0(C, \omega_C^{\otimes 2}(D))$, $u \mapsto B_E(v,u)$ has rank $r$. Then $\on{rk}(E)\geq g-r.$"
  \textit{Source: Landesman--Litt, §5.2 (proposition:pairing-result).}
  Let $C$ be a smooth proper connected curve of genus $g$, $D \subset C$ a
  reduced divisor, and $B_E\colon (E \otimes \omega_C(D)) \times
  (E^\vee \otimes \omega_C) \to \omega_C^{\otimes 2}(D)$ the natural perfect
  pairing obtained by tensoring and applying the trace pairing
  $E \otimes E^\vee \to \mathcal{O}_C$.
  Let $E_\star$ be a semistable parabolic bundle on $(C, D)$ of parabolic
  degree zero with underlying vector bundle $E$.
  Let $v \in H^0(C, E \otimes \omega_C(D))$ be a nonzero global section, and
  suppose the induced map
  \[
    H^0(C, B_E(v, -))\colon H^0(C, E^\vee \otimes \omega_C)
    \to H^0(C, \omega_C^{\otimes 2}(D)),
    \qquad u \mapsto B_E(v, u),
  \]
  has rank $r$.
  Then $\rk(E) \geq g - r$.
\end{proposition}

\begin{proof}
  \uses{thm:period-map}
  % SOURCE QUOTE PROOF: "The idea is to apply \autoref{proposition:generic-parabolic-global-generation}(II) with $c =1$ and $\delta = r$, and we now set up notation to do so. Set $n=\deg D$. Let $f_v: E^\vee\otimes \omega_C\to \omega_C^{\otimes 2}(D)$ be the map of vector bundles induced by $B_E(v, -)$, and set $U=\ker(f_v)$. As $B_E$ is a perfect pairing and $v$ is nonzero by assumption, $U$ has corank one in $E^\vee\otimes \omega_C$. We conclude by applying \autoref{proposition:generic-parabolic-global-generation}(II) to the parabolic bundle $(E_\star)^\vee \otimes \omega_C(D)$."
  Let $f_v\colon E^\vee \otimes \omega_C \to \omega_C^{\otimes 2}(D)$ be the map
  of vector bundles induced by $B_E(v, -)$, and set $U = \ker(f_v)$.
  Since $B_E$ is a perfect pairing and $v$ is nonzero, $U$ has corank one in
  $E^\vee \otimes \omega_C$.
  Applying the generic global generation result for semistable parabolic
  bundles (Landesman--Litt--Sawin, [LL:geometric-local-systems] Prop.~6.3.6),
  case~(II), with $c = 1$ and $\delta = r$, to the Serre-dual parabolic bundle
  $(E_\star)^\vee \otimes \omega_C(D)$ yields $\rk(E) \geq g - r$.
\end{proof}

\section{The main cohomological results}

% ---------------------------------------------------------------------------
% Rank bound (§6.1) and vanishing result (§6.2)
% ---------------------------------------------------------------------------

\begin{theorem}[Cohomological rank bound]
  \label{thm:rank-bound-subsystem}
  \lean{MR4717077CanonicalRepresentationsOfSurfaceGroups.rankBound_subsystem}
  \uses{thm:period-map, proposition:pairing-result}
  % SOURCE: [Landesman--Litt], Theorem 1.7.1 (theorem:rank-bound-subsystem), (read from references/MR4717077-surface-groups.tex)
  % SOURCE QUOTE: "Let $\pi:\mathscr{C}\to \mathscr{M}$ be a smooth proper family of $n$-pointed curves of genus $g$ with geometrically connected fibers, so that the associated map $\mathscr{M}\to \mathscr{M}_{g,n}$ is dominant \'etale. Let $\pi^\circ: \mathscr{C}^\circ\to \mathscr{M}$ be the associated family of punctured curves. Let $\mathbb{V}$ be a unitary local system on $\mathscr{C}^\circ$. Then any nonzero sub-local system of $R^1\pi_*^\circ \mathbb{V}$ has rank at least $2g-2\rk\mathbb{V}$."
  \textit{Source: Landesman--Litt, Theorem 1.7.1.}
  Let $\pi\colon \mathcal{C} \to \mathcal{M}$ be a smooth proper family of
  $n$-pointed genus-$g$ curves with geometrically connected fibers such that
  $\mathcal{M} \to \mathcal{M}_{g,n}$ is dominant \'{e}tale.
  Let $\mathbb{V}$ be a unitary local system on the associated family of
  punctured curves $\mathcal{C}^\circ$.
  Then any nonzero sub-local system of $R^1\pi^\circ_* \mathbb{V}$ has rank
  at least $2g - 2\rk \mathbb{V}$.
  In particular, $H^0(\mathcal{M},\, R^1\pi^\circ_* \mathbb{V}) = 0$ whenever
  $g > \rk \mathbb{V}$.
\end{theorem}

\begin{proof}
  \uses{thm:period-map, proposition:pairing-result}
  A sub-local system of $R^1\pi^\circ_* \mathbb{V}$ of rank below
  $2g - 2\rk\mathbb{V}$ would, by the theorem of the fixed part, force the
  derivative of the period map to have a kernel.
  By \cref{thm:period-map}, that derivative is the multiplication map on
  global sections induced by the bilinear pairing $B_E$ on
  $E \otimes \omega_C(D)$ versus $E^\vee \otimes \omega_C$.
  Such a kernel produces a nonzero $v$ for which the rank of $B_E(v, -)$ is
  small; \cref{proposition:pairing-result} then forces $\rk(E) = \rk\mathbb{V}$
  to be large, yielding the required rank bound.
\end{proof}

\begin{theorem}[Cohomological vanishing]
  \label{thm:cohomological-vanishing}
  \lean{MR4717077CanonicalRepresentationsOfSurfaceGroups.cohomologicalVanishing}
  \uses{thm:rank-bound-subsystem, lem:tensor-product-unitary}
  % SOURCE: [Landesman--Litt], Theorem 6.2.1 (theorem:unitary-rigid), (read from references/MR4717077-surface-groups.tex)
  % SOURCE QUOTE: "With notation as in \autoref{notation:versal-family}, let $\pi^\circ: \mathscr{C}^\circ\to \mathscr{M}$ be a punctured versal family of genus $g$ curves. Let $m\in \mathscr{M}$ be a point and set $C^\circ = \mathscr C^\circ_m$. Let $A$ be a local Artin $\mathbb{C}$-algebra with residue field $\mathbb{C}$, and let $\mathbb V$ be a locally constant sheaf of free $A$-modules on $\mathscr C^\circ$ such that (1) $\mathbb V|_{C^\circ}$ is a constant deformation of a unitary local system, i.e. there exists a unitary complex local system $\mathbb{V}_0$ on $C^\circ$ and an isomorphism $\mathbb V|_{C^\circ}\simeq \mathbb{V}_0\otimes A$, and (2) $\rk_A \mathbb V < g$. Then $H^0(\mathscr M, R^1 \pi^\circ_* \mathbb V) = 0$."
  \textit{Source: Landesman--Litt, Theorem 6.2.1.}
  Let $\pi^\circ\colon \mathcal{C}^\circ \to \mathcal{M}$ be a punctured versal
  family of genus-$g$ curves.
  Let $\mathbb{V}$ be a local system on $\mathcal{C}^\circ$ whose restriction
  to a fiber $C^\circ$ is unitary, with $\rk \mathbb{V} < g$.
  Then $H^0(\mathcal{M},\, R^1\pi^\circ_* \mathbb{V}) = 0$.
  In particular, $R^1\pi^\circ_* \mathbb{V}$ has no nonzero unitary
  sub-local systems.
\end{theorem}

\begin{proof}
  \uses{thm:rank-bound-subsystem, lem:tensor-product-unitary}
  By \cref{lem:tensor-product-unitary}, after a dominant \'{e}tale base change
  write $\mathbb{V} \cong \bigoplus_i \mathbb{U}_i \otimes (\pi^\circ)^* \mathbb{W}_i$
  with $\mathbb{U}_i$ unitary.
  A nonzero global section of $R^1\pi^\circ_*(\mathbb{U}_i \otimes (\pi^\circ)^* \mathbb{W}_i)
  = (R^1\pi^\circ_* \mathbb{U}_i) \otimes \mathbb{W}_i$ yields a nonzero map
  $\mathbb{W}_i^\vee \to R^1\pi^\circ_* \mathbb{U}_i$, i.e.\ a sub-local system
  of rank $\rk\mathbb{W}_i$.
  Since $\mathbb{U}_i$ is unitary, \cref{thm:rank-bound-subsystem} gives
  $\rk\mathbb{W}_i \geq 2g - 2\rk\mathbb{U}_i$, so
  $\rk\mathbb{V} \geq \rk\mathbb{U}_i \cdot \rk\mathbb{W}_i \geq g$,
  contradicting $\rk\mathbb{V} < g$.
\end{proof}

\section{The asymptotic Putman-Wieland conjecture}

% ---------------------------------------------------------------------------
% Putman--Wieland (§7.2 Our results)
% ---------------------------------------------------------------------------

\begin{theorem}[Asymptotic Putman--Wieland]
  \label{thm:putman-wieland-asymptotic}
  \lean{MR4717077CanonicalRepresentationsOfSurfaceGroups.putmanWieland_asymptotic}
  \uses{thm:rank-bound-subsystem, lem:tensor-product-unitary}
  % SOURCE: [Landesman--Litt], Theorem 7.2.1 (theorem:asymptotic-putman-wieland), (read from references/MR4717077-surface-groups.tex)
  % SOURCE QUOTE: "With notation as in \autoref{notation:prym}, let $\rho$ be an irreducible complex $H$-representation and $\Gamma' \subset \Gamma$ be a finite-index subgroup. Then $H^1(\Sigma_{g',b',p'},\mathbb C)^\rho$ has no nonzero $\Gamma'$-invariant subrepresentations of dimension strictly less than $2g - 2\dim \rho$. The same holds for $H_1(\Sigma_{g'},\mathbb C)^\rho$ in place of $H^1(\Sigma_{g',b',p'},\mathbb C)^\rho$."
  \textit{Source: Landesman--Litt, Theorem 7.2.1.}
  With notation as in the higher Prym setup
  (finite Galois cover $\Sigma_{g',b',p'} \to \Sigma_{g,b,p}$ with Galois group
  $H$, corresponding to a surjection
  $\phi\colon \pi_1(\Sigma_{g,b,p}) \twoheadrightarrow H$ with kernel $K$, and
  $\Gamma \subset \PMod_{g,b,p+1}$ the finite-index subgroup preserving $\phi$
  up to conjugacy), let $\rho$ be an irreducible complex $H$-representation and
  $\Gamma' \subset \Gamma$ a finite-index subgroup.
  Then the $\rho$-isotypic component $H^1(\Sigma_{g',b',p'}, \mathbb{C})^\rho$
  has no nonzero $\Gamma'$-invariant subrepresentation of dimension strictly
  less than $2g - 2\dim\rho$.
  The same holds for $H_1(\Sigma_{g'}, \mathbb{C})^\rho$ in place of
  $H^1(\Sigma_{g',b',p'}, \mathbb{C})^\rho$.
\end{theorem}

\begin{proof}
  \uses{thm:rank-bound-subsystem, lem:tensor-product-unitary}
  The $H$-cover produces finite-monodromy (hence unitary) local systems
  $\mathbb{U}_\rho$ on $\mathcal{C}^\circ$ with monodromy given by the
  irreducible representations of $H$.
  After passing to a finite \'{e}tale cover to trivialize $\mathbb{W}_\rho$
  (via \cref{lem:tensor-product-unitary}), a $\Gamma$-stable subspace of
  $H_1(\Sigma_{g'})^\rho$ of dimension less than $2g - 2\dim\rho$ would yield
  a sub-local system of $R^1\pi^\circ_* \mathbb{U}_\rho$ of rank below
  $2g - 2\rk\mathbb{U}_\rho$, contradicting \cref{thm:rank-bound-subsystem}.
\end{proof}

\begin{corollary}[Putman--Wieland for small $|H|$]
  \label{cor:putman-wieland-small-H}
  \lean{MR4717077CanonicalRepresentationsOfSurfaceGroups.putmanWieland_smallH}
  \uses{thm:putman-wieland-asymptotic}
  % SOURCE: [Landesman--Litt], Corollary 7.2.4 (corollary:asymptotic-putman-wieland), (read from references/MR4717077-surface-groups.tex)
  % SOURCE QUOTE: "For fixed $g, b, p \geq 0$, $\PW_{g,b,p}^H$ holds for every group $H$ with $\# H < g^2$."
  \textit{Source: Landesman--Litt, Corollary 7.2.4.}
  For fixed $g, b, p \geq 0$ and any finite group $H$ with $\#H < g^2$, the
  statement $\PW_{g,b,p}^H$ holds: for any surjection
  $\phi\colon \pi_1(\Sigma_{g,b,p}) \twoheadrightarrow H$ with kernel $K$, every
  nonzero vector in $V_K = H_1(\Sigma_{g'}, \mathbb{C})$ has infinite orbit
  under $\Gamma$.
\end{corollary}

\begin{proof}
  \uses{thm:putman-wieland-asymptotic}
  Since $\#H < g^2$ and $\#H$ equals the sum of squares of dimensions of
  irreducible representations of $H$, every irreducible $\rho$ of $H$
  satisfies $\dim\rho < g$.
  Apply \cref{thm:putman-wieland-asymptotic}: any $\Gamma$-stable subspace of
  $H_1(\Sigma_{g'})^\rho$ has dimension at least $2g - 2\dim\rho > 0$, so
  no nonzero invariant vector exists.
\end{proof}

\section{Proof of the main theorem on MCG-finite representations}

% ---------------------------------------------------------------------------
% Background on cohomological rigidity (§8.1) and integrality
% ---------------------------------------------------------------------------

\begin{theorem}[Cohomologically rigid implies integral]
  \label{thm:cohomologically-rigid-integral}
  \lean{MR4717077CanonicalRepresentationsOfSurfaceGroups.cohomRigid_isIntegral}
  An irreducible cohomologically rigid local system on a smooth complex variety
  with quasi-unipotent local monodromy around boundary components of a good
  compactification is defined over the ring of integers $\mathcal{O}_K$ of some
  number field $K$.
  (Axiom: main result of Esnault--Groechenig \cite{esnault2018cohomologically}
  and Klevdal--Patrikis \cite{klevdalP:g-rigid-local-systems-are-integral};
  not yet available in Mathlib.)
\end{theorem}

\begin{proof}
  Proved by Esnault--Groechenig (arXiv:1711.06427) and, in the generality
  needed here, by Klevdal--Patrikis.  The argument combines $p$-adic Hodge
  theory and de Rham cohomology to extract an integral structure from
  cohomological rigidity.  Introduced as an axiom in the Lean formalization;
  the relevant $p$-adic Hodge theory infrastructure is not yet in Mathlib.
\end{proof}

\begin{proposition}[Isomonodromy rigidity]
  \label{lem:isomonodromy-rigidity}
  \lean{MR4717077CanonicalRepresentationsOfSurfaceGroups.isomonodromy_rigidity}
  \uses{thm:cohomologically-rigid-integral}
  % SOURCE: [Landesman--Litt], Proposition 8.2.1 (proposition:cohomologically-rigid), (read from references/MR4717077-surface-groups.tex)
  % SOURCE QUOTE: "Let $\mathbb{V}$ be a $\on{GL}_r$ (respectively, ~$\on{PGL}_r$)-local system on $\mathscr{C}^\circ$ with $r<\sqrt{g+1}$. Suppose that for $m\in \mathscr{M}$, $\mathbb{V}|_{{C}^\circ}$ is (respectively, is the projectivization of) an irreducible, unitary local system. Then $\mathbb{V}$ is strongly cohomologically rigid. In particular, $\mathbb V$ is cohomologically rigid."
  \textit{Source: Landesman--Litt, Proposition 8.2.1.}
  Let $\mathbb{V}$ be an irreducible $\GL_r$-local system on $\mathcal{C}^\circ$
  with $r < \sqrt{g+1}$, whose restriction to a fiber $C^\circ$ is unitary.
  Then $\mathbb{V}$ is strongly cohomologically rigid; in particular, it is
  cohomologically rigid and its isomonodromic deformations are algebraic.
\end{proposition}

\begin{proof}
  \uses{thm:cohomological-vanishing, thm:cohomologically-rigid-integral}
  Strong cohomological rigidity requires $H^1(\mathcal{C}^\circ, \mathrm{ad}\,\mathbb{V}) = 0$.
  By the Leray spectral sequence, it suffices to show
  $H^0(\mathcal{M}, R^1\pi^\circ_* \mathrm{ad}\,\mathbb{V}) = 0$ and
  $H^1(\mathcal{M}, \pi^\circ_* \mathrm{ad}\,\mathbb{V}) = 0$.
  The first vanishes by \cref{thm:cohomological-vanishing} since
  $\rk(\mathrm{ad}\,\mathbb{V}) = r^2-1 < g$ (because $r < \sqrt{g+1}$) and
  $\mathrm{ad}\,\mathbb{V}|_{C^\circ}$ is unitary.
  The second vanishes by Schur's lemma since $\mathbb{V}|_{C^\circ}$ is
  irreducible.
  By \cref{thm:cohomologically-rigid-integral}, the cohomological rigidity
  implies $\mathbb{V}$ is defined over $\mathcal{O}_K$, so isomonodromic
  deformations are algebraic.
\end{proof}

% ---------------------------------------------------------------------------
% Unitary case (§8.4), semisimple case (§8.5), semisimplicity (§8.6),
% completion of the main theorem (§8)
% ---------------------------------------------------------------------------

\begin{proposition}[Unitary case]
  \label{prop:unitary-case}
  \lean{MR4717077CanonicalRepresentationsOfSurfaceGroups.finiteImage_unitary}
  \uses{thm:cohomologically-rigid-integral, thm:cohomological-vanishing, thm:nah-unitarity}
  % SOURCE: [Landesman--Litt], Proposition 8.4.1 (proposition:unitary-implies-finite), (read from references/MR4717077-surface-groups.tex)
  % SOURCE QUOTE: "Let  $$\rho:\pi_1(\Sigma_{g,n})\to \on{GL}_r(\mathbb{C})$$ be a unitary MCG-finite representation with $r<\sqrt{g+1}$. Then $\rho$ has finite image."
  \textit{Source: Landesman--Litt, Proposition 8.4.1.}
  Let $\rho\colon \pi_1(\Sigma_{g,n}) \to \GL_r(\mathbb{C})$ be a unitary
  MCG-finite representation with $r < \sqrt{g+1}$.
  Then $\rho$ has finite image.
\end{proposition}

\begin{proof}
  \uses{lem:isomonodromy-rigidity, thm:cohomologically-rigid-integral, prop:mcg-finite-basic, thm:nah-unitarity}
  By \cref{prop:mcg-finite-basic} it suffices to treat the irreducible case.
  By \cref{lem:isomonodromy-rigidity}, $\rho$ is cohomologically rigid.
  By \cref{thm:cohomologically-rigid-integral} (Klevdal--Patrikis), $\rho$ is
  defined over $\mathcal{O}_K$ for some number field $K$.
  For each embedding $\tau\colon \mathcal{O}_K \hookrightarrow \mathbb{C}$,
  the twist $\rho_\tau = \rho \otimes_{\mathcal{O}_K,\tau} \mathbb{C}$ is also
  cohomologically rigid, irreducible, and MCG-finite; by
  \cref{thm:nah-unitarity} each $\rho_\tau$ is unitary.
  Since $\rho$ takes values in $\GL_r(\mathcal{O}_K)$ with every embedding
  giving a unitary representation, compactness of $U(r)$ and discreteness of
  $\mathcal{O}_K$ give finite image.
\end{proof}

\begin{theorem}[Semisimple case]
  \label{thm:semisimple-case}
  \lean{MR4717077CanonicalRepresentationsOfSurfaceGroups.finiteImage_semisimple}
  \uses{thm:cohomological-vanishing, prop:unitary-case, lem:isomonodromy-rigidity, prop:mcg-finite-basic, prop:mcg-finite-family, thm:nah-unitarity}
  % SOURCE: [Landesman--Litt], Theorem 8.5.3 (theorem:finite-image-semisimple), (read from references/MR4717077-surface-groups.tex)
  % SOURCE QUOTE: "Let $$\rho: \pi_1(\Sigma_{g,n})\to \on{GL}_r(\mathbb{C})$$ be a semisimple MCG-finite representation, and suppose $r<\sqrt{g+1}$. Then $\rho$ has finite image."
  \textit{Source: Landesman--Litt, Theorem 8.5.3.}
  Let $\rho\colon \pi_1(\Sigma_{g,n}) \to \GL_r(\mathbb{C})$ be a semisimple
  MCG-finite representation with $r < \sqrt{g+1}$.
  Then $\rho$ has finite image.
\end{theorem}

\begin{proof}
  \uses{prop:mcg-finite-basic, prop:mcg-finite-family, prop:unitary-case, thm:cohomological-vanishing}
  By \cref{prop:mcg-finite-basic} it suffices to treat the irreducible case.
  Lift $\rho$ to a local system $\mathbb{V}$ on a punctured versal family
  $\pi^\circ\colon \mathcal{C}^\circ \to \mathcal{M}$ (converse of
  \cref{prop:mcg-finite-family}).
  By non-abelian Hodge theory, deform $\mathbb{V}$ (with constant determinant)
  to $\mathbb{V}_0$ underlying a complex polarizable variation of Hodge
  structure.
  Since $r < 2\sqrt{g+1}$, Landesman--Litt [LL22] Theorem~1.2.12 gives that
  $\mathbb{V}_0|_{C^\circ}$ is unitary; it is MCG-finite by
  \cref{prop:mcg-finite-family}; hence by \cref{prop:unitary-case} it has
  finite image.
  Finally, $\mathbb{V}|_{C^\circ} = \mathbb{V}_0|_{C^\circ}$: the unitary
  fibral monodromy of $\mathbb{V}_0$ admits no nontrivial MCG-finite
  deformations by \cref{thm:cohomological-vanishing} (applied as in
  Lemma~8.3.2 of the paper), so the deformation was trivial on fibers.
\end{proof}

\begin{lemma}[Semisimplicity of low-rank MCG-finite representations]
  \label{lem:semisimplicity}
  \lean{MR4717077CanonicalRepresentationsOfSurfaceGroups.mcgFinite_isSemisimple}
  \uses{thm:putman-wieland-asymptotic}
  % SOURCE: [Landesman--Litt], Lemma 8.6.1 (lemma:semisimplicity), (read from references/MR4717077-surface-groups.tex)
  % SOURCE QUOTE: "Let $\rho$ be an MCG-finite representation of $\pi_1(\Sigma_{g,n})$, and let $\rho_1\subset\rho$ be a semisimple, characteristic subrepresentation of $\rho$. Let $\rho_2=\rho/\rho_1$, and suppose $\rho_2$ is semisimple as well. If $\rho_1, \rho_2$ have finite image and $\rk\rho<\sqrt{g+1}$, then $\rho$ has finite image as well, i.e.~the extension of $\rho_2$ by $\rho_1$ splits."
  \textit{Source: Landesman--Litt, Lemma 8.6.1.}
  Let $\rho$ be an MCG-finite representation of $\pi_1(\Sigma_{g,n})$ with
  $\rk\rho < \sqrt{g+1}$.
  Let $\rho_1 \subset \rho$ be a semisimple characteristic subrepresentation,
  $\rho_2 = \rho/\rho_1$ semisimple, with $\rho_1, \rho_2$ both of finite
  image.
  Then $\rho$ has finite image, i.e.\ the extension of $\rho_2$ by $\rho_1$
  splits.
  In particular, every MCG-finite representation of $\pi_1(\Sigma_{g,n})$
  with $\rk < \sqrt{g+1}$ is semisimple.
\end{lemma}

\begin{proof}
  \uses{thm:putman-wieland-asymptotic}
  Pass to a finite Galois cover $\Sigma_{g',n'} \to \Sigma_{g,n}$ (Galois
  group $H$) trivializing $\rho_1$ and $\rho_2$.
  The extension class lies in
  $\Ext^1_{\pi_1(\Sigma_{g',n'})}(\rho_2, \rho_1)
  = \mathrm{Hom}(\pi_1(\Sigma_{g',n'}), \mathrm{Hom}(\rho_2,\rho_1))$,
  which factors through $H_1(\Sigma_{g',n'})$.
  Decompose $\rho_2^\vee \otimes \rho_1 = \bigoplus_i \sigma_i^{n_i}$
  irreducibly; since $\rk\rho < \sqrt{g+1}$, the AM-GM inequality gives
  $n_i \dim\sigma_i < (g+1)/4$, so $\dim\sigma_i < g$.
  The projection of the extension class onto the $\sigma_i$-isotypic part of
  $H_1(\Sigma_{g',n'})$ gives a $\Gamma$-stable subspace of rank at most
  $(g+1)/4$.
  But \cref{thm:putman-wieland-asymptotic} shows any nonzero such subspace has
  rank at least $2g - 2\dim\sigma_i > (g+1)/4$, so each projection is zero
  and $\rho$ has finite image.
\end{proof}

\begin{theorem}[Finite image of MCG-finite representations]
  \label{thm:finite-image}
  \lean{MR4717077CanonicalRepresentationsOfSurfaceGroups.finiteImage_of_mcgFinite}
  \uses{def:mcg-finite, thm:semisimple-case, lem:semisimplicity, prop:mcg-finite-basic}
  % SOURCE: [Landesman--Litt], Theorem 1.2.1 (theorem:finite-image), (read from references/MR4717077-surface-groups.tex)
  % SOURCE QUOTE: "For $g, n, r\geq 0$, let $$\rho: \pi_1(\Sigma_{g,n})\to \on{GL}_r(\mathbb{C})$$ be a MCG-finite representation. If $r<\sqrt{g+1}$, then $\rho$ has finite image."
  \textit{Source: Landesman--Litt, Theorem 1.2.1.}
  For $g, n, r \geq 0$, let
  $\rho\colon \pi_1(\Sigma_{g,n}) \to \GL_r(\mathbb{C})$ be a MCG-finite
  representation.
  If $r < \sqrt{g+1}$, then $\rho$ has finite image.
\end{theorem}

\begin{proof}
  \uses{thm:semisimple-case, lem:semisimplicity, prop:mcg-finite-basic}
  By induction on the Jordan--H\"older length of $\rho$.
  By \cref{prop:mcg-finite-basic}, each Jordan--H\"older factor is MCG-finite
  with rank less than $\sqrt{g+1}$, hence has finite image by
  \cref{thm:semisimple-case} applied inductively.
  \cref{lem:semisimplicity} then shows that successive extensions split, so
  $\rho$ has finite image.
\end{proof}

\section{Consequences for arithmetic representations}

% ---------------------------------------------------------------------------
% Relative Fontaine--Mazur (§9.1) and lifting residual representations (§9.2)
% ---------------------------------------------------------------------------

\begin{theorem}[Arithmetic consequence: relative Fontaine--Mazur]
  \label{theorem:arithmetic-consequence}
  \lean{MR4717077CanonicalRepresentationsOfSurfaceGroups.arithmetic_finiteImage}
  \uses{thm:finite-image}
  % SOURCE: [Landesman--Litt], §9.1 (theorem:arithmetic-consequence), (read from references/MR4717077-surface-groups.tex)
  % SOURCE QUOTE: "Let $K$ be a finitely-generated field of characteristic zero and $(C, x_1, \cdots, x_n)$ a smooth $n$-pointed geometrically connected curve of genus $g$ over $K$, such that the corresponding map $\on{Spec}(K)\to \mathscr{M}_{g,n, \mathbb{Q}}$ factors through the generic point. If $r<\sqrt{g+1}$, any continuous arithmetic representation $\rho: \pi_1^{\text{\'et}}(C_{\overline{K}}\setminus\{x_1, \cdots, x_n\})\to \on{GL}_r(\overline{\mathbb{Q}}_\ell)$ has finite image."
  \textit{Source: Landesman--Litt, §9.1 (theorem:arithmetic-consequence).}
  Let $K$ be a finitely-generated field of characteristic zero and
  $(C, x_1, \ldots, x_n)$ a smooth $n$-pointed geometrically connected curve of
  genus $g$ over $K$ such that the corresponding map
  $\mathrm{Spec}(K) \to \mathcal{M}_{g,n,\mathbb{Q}}$ factors through the
  generic point.
  If $r < \sqrt{g+1}$, then any continuous arithmetic representation
  \[
    \rho\colon \pi_1^{\text{\'et}}(C_{\overline{K}} \setminus \{x_1, \ldots, x_n\})
    \to \GL_r(\overline{\mathbb{Q}}_\ell)
  \]
  has finite image.
  (Here $\rho$ is \emph{arithmetic} if its conjugacy class has finite orbit
  under the action of $\mathrm{Gal}(\overline{K}/K)$ on isomorphism classes of
  $\pi_1^{\text{\'et}}(C_{\overline{K}} \setminus \{x_1, \ldots, x_n\})$-representations.)
\end{theorem}

\begin{proof}
  \uses{thm:finite-image, prop:mcg-finite-family}
  % SOURCE QUOTE PROOF: "The two exact sequences in (map-of-birman-exact-sequences) induce compatible outer actions of $\pi_1(\on{Spec} K)$ and $\pi_1(\mathscr M_{g,n,\mathbb Q})$... Because $\on{Spec} K \to \mathscr M_{g,n,\mathbb Q}$ is dominant, the induced map $\pi_1(\on{Spec} K) \to \pi_1(\mathscr M_{g,n,\mathbb Q})$ has image of finite index... Hence $\rho|_{\pi_1((\mathscr{C}^\circ)^{\text{an}})}$ is MCG-finite, and hence has finite image by \autoref{theorem:finite-image}, which may be applied after choosing an isomorphism between $\overline{\mathbb{Q}}_\ell$ and $\mathbb{C}$. But $\pi_1( (\mathscr{C}^\circ)^{\text{an}})$ is dense in $\pi_1^{\text{\'et}}(\mathscr{C}^\circ_{\overline{K}})$, so $\rho$ also has finite image."
  Comparing the homotopy exact sequence of $C^\circ_{\overline{K}} \to
  \mathrm{Spec}\,K$ with the Birman-type sequence for the universal punctured
  curve over $\mathcal{M}_{g,n,\mathbb{Q}}$, the outer action of
  $\pi_1(\mathrm{Spec}\,K)$ on $\pi_1^{\text{\'et}}(C^\circ_{\overline{K}})$
  factors through $\pi_1^{\text{\'et}}(\mathcal{M}_{g,n,\mathbb{Q}})$.
  Since $\mathrm{Spec}\,K \to \mathcal{M}_{g,n,\mathbb{Q}}$ is dominant, this map
  has image of finite index, so the conjugacy class of $\rho$ has finite orbit
  under $\pi_1^{\text{\'et}}(\mathcal{M}_{g,n,\overline{\mathbb{Q}}})$.
  Choosing an embedding $\overline{K} \hookrightarrow \mathbb{C}$ and an
  isomorphism $\mathbb{C} \cong \overline{\mathbb{Q}}_\ell$, the standard
  comparison theorems show $\rho$ restricted to the topological fundamental
  group of the associated Riemann surface is MCG-finite; hence it has finite
  image by \cref{thm:finite-image}.
  Since the topological fundamental group is dense in the \'etale one, $\rho$
  itself has finite image.
\end{proof}

\begin{corollary}[No arithmetic lifts of residual representations]
  \label{corollary:lifting-and-geometric-origin}
  \lean{MR4717077CanonicalRepresentationsOfSurfaceGroups.noArithmeticLift}
  \uses{theorem:arithmetic-consequence}
  % SOURCE: [Landesman--Litt], §9.2 (corollary:lifting-and-geometric-origin), (read from references/MR4717077-surface-groups.tex)
  % SOURCE QUOTE: "Let $K, (C, x_1, \cdots, x_n)$ be as in \autoref{theorem:arithmetic-consequence}. Then for $1<r<\sqrt{g+1}$ and $p\gg_r 0$, no surjective representation $\rho: \pi_1^{\text{\'et}}(C_{\overline{K}}\setminus\{x_1, \cdots, x_n\})\to \on{GL}_r(\mathbb{F}_p)$ admits an arithmetic lift to characteristic zero. In particular no such representation is of geometric origin."
  \textit{Source: Landesman--Litt, §9.2 (corollary:lifting-and-geometric-origin).}
  Let $K, (C, x_1, \ldots, x_n)$ be as in
  \cref{theorem:arithmetic-consequence}.
  Then for $1 < r < \sqrt{g+1}$ and $p \gg_r 0$, no surjective representation
  \[
    \rho\colon \pi_1^{\text{\'et}}(C_{\overline{K}} \setminus \{x_1, \ldots, x_n\})
    \to \GL_r(\mathbb{F}_p)
  \]
  admits an arithmetic lift to characteristic zero.
  In particular, no such representation is of geometric origin.
\end{corollary}

\begin{proof}
  \uses{theorem:arithmetic-consequence}
  % SOURCE QUOTE PROOF: "Suppose to the contrary that it did admit an arithmetic lift $\xi$; by \autoref{theorem:arithmetic-consequence}, $\xi$ would have finite image. Thus it suffices to show that for $Q$ a finite, totally ramified extension of $\mathbb{Q}_p$, there do not exist any finite subgroups of $\on{GL}_r(\mathscr{O}_Q)$ surjecting onto $\on{GL}_r(\mathbb{F}_p)$ if $p\gg_r 0$. This follows from Jordan's theorem on finite subgroups of $\on{GL}_r(Q)$..."
  A residual representation of geometric origin lifts to one of geometric
  origin over $\overline{\mathbb{Q}}_p$, and such lifts are arithmetic, so it
  suffices to rule out arithmetic lifts.
  By \cref{theorem:arithmetic-consequence}, any arithmetic lift $\xi$ would have
  finite image, hence land in a finite subgroup of $\GL_r(\mathcal{O}_Q)$ for
  some finite totally ramified extension $Q/\mathbb{Q}_p$.
  By Jordan's theorem, such a finite subgroup contains an abelian normal
  subgroup of index dividing a constant $n(r)$; explicit unipotent generators
  $\overline{g_1}, \overline{g_2} \in \GL_r(\mathbb{F}_p)$ violate the resulting
  commutation relation whenever $p \nmid n(r)$.
  Thus for $p \gg_r 0$ no finite subgroup of $\GL_r(\mathcal{O}_Q)$ surjects
  onto $\GL_r(\mathbb{F}_p)$, so no such $\rho$ admits an arithmetic lift.
\end{proof}
